% !PRINTED 757
\weberpage{757}

Nehmen wir also an, es sei $\pi$ und folglich auch $i\pi$ eine algebraische Zahl, so ist $i\pi$ eine der Wurzeln einer irreduciblen rationalen Gleichung $\chi(x)=0$, deren Coëfficienten ganze rationale Zahlen sind.

Bezeichnen wir die sämmtlichen Wurzeln dieser algebraischen Gleichung mit $\beta_1,\beta_2,\ldots,\beta_\nu$, und bezeichnen den Coëfficienten von $x^\nu$ in $\chi$ mit $a$, so ist
\weq{2}{\chi(x) = a(x-\beta_1)(x-\beta_2)\ldots(x-\beta_\nu),}
und die Producte $a\beta_1,a\beta_2,\ldots,a\beta_\nu$ sind \gesperrt{ganze} algebraische Zahlen, ihre ganzen symmetrischen Functionen also \gesperrt{ganze rationale Zahlen} (\S.\,133).

Da nun die Zahl $i\pi$ unter den $\beta$ vorkommen soll, so ist nach (1):
\[
(1+e^{\beta_1})(1+e^{\beta_2})\ldots(1+e^{\beta_\nu}) = 0,
\]
und wenn wir die Multiplication ausführen, so ergiebt sich eine Gleichung:
\[
1 + \sum e^{\beta_i} + \sum e^{\beta_i+\beta_h} + \sum e^{\beta_i+\beta_h+\beta_k} + \cdots = 0.
\]
Unter den Exponenten in dieser Summe kann mehrmals die Zahl Null vorkommen; wir wollen annehmen $(C-1)$mal, so dass $C$ eine positive ganze Zahl, mindestens $=1$, ist. Die übrigen Exponenten $\beta_i,\beta_i+\beta_h,\beta_i+\beta_h+\beta_k,\ldots$, die zum Theil auch unter einander gleich sein können, wollen wir mit $\alpha_1,\alpha_2,\ldots,\alpha_\mu$ bezeichnen, so dass die Gleichung besteht:
\weq{3}{C + e^{\alpha_1} + e^{\alpha_2} + \cdots + e^{\alpha_\mu} = 0.}
Hierin sind nun die $\alpha_1,\alpha_2,\ldots,\alpha_\mu$ algebraische Zahlen, die, mit der ganzen rationalen Zahl $a$ multiplicirt, in ganze algebraische Zahlen übergehen.

Die symmetrischen Functionen der sämmtlichen Summen $\beta_i,\beta_i+\beta_h,\beta_i+\beta_h+\beta_k,\ldots$ sind aber zugleich symmetrische Functionen der $\beta_1,\ldots,\beta_\nu$, und folglich rationale Zahlen. Diese Summen sind folglich auch die Wurzeln einer rationalen Gleichung, und da man die Wurzel $0$, so oft sie vorhanden ist, absondern kann, so sind auch die $\alpha_1,\alpha_2,\ldots,\alpha_\mu$ die Wurzeln einer rationalen Gleichung; die symmetrischen Grundfunctionen der $a\alpha_1,a\alpha_2,\ldots,a\alpha_\mu$ sind \gesperrt{ganze rationale Zahlen}.

Die absoluten Werthe der Zahlen
\[
\alpha_1,\alpha_2,\ldots,\alpha_\mu
\]
sollen jetzt mit
\[
a_1,a_2,\ldots,a_\mu
\]
bezeichnet werden.
